\subsection{首选最小实例簇的组装：从 LOST–Fold–Čech–prime–budget 到最小可审计时空实例}\label{subsec:typed-address-biaxial-completion-preferred-minimal-instantiation}
上一节已经说明，若想把本文的统一骨架真正实例化为物理型时空候选，则对象层至少要同时保留六类结构：状态—forcing—更新、跨分辨率逆系统、局部覆盖与截面、谱接口、残差外置以及资源预算。
这一要求并不是额外附会，而是由正文各处已经建立的对象共同强迫出来的：观察者时空逻辑系统给出状态、事件、区域与通信；Fold 的 restriction 逆系统给出跨分辨率的一致本体；\(\mathrm{NULL}\) 三分解与 Čech 型粘合障碍给出局部—全局失败模式；prime-register 给出乘法层与前史层的不可有限化账本；而高阶碰撞矩、压力、\(\zeta\) 与 projective transfer 又把可见统计、primitive 原子与机制谱焊成同一谱层。
这意味着，本文事实上已经隐含拥有一个首选最小实例簇，只差把它明确组装出来。

本节的目标正是完成这一步：给出一组最小而完整的数据，使得上一节的粗糙时空四元组
\[
\mathfrak G=(\preceq,\tau,d_{\mathrm{res}},\Omega)
\]
不再只是抽象模板，而成为一个可以逐项审计、逐项替换、逐项数值化的具体实例。

\subsubsection{首选最小实例簇的定义}\label{subsubsec:typed-address-biaxial-completion-preferred-minimal-instantiation-definition}
先把六元组具体写出。

\begin{definition}[首选最小实例簇]\label{def:typed-address-biaxial-completion-preferred-minimal-instantiation}
定义
\[
\mathcal I_{\min}
:=
\bigl(
\mathcal P_{\mathrm{LOST}},
\mathcal X_{\mathrm{Fold}},
\mathcal U_{\mathrm{loc}},
\mathcal R_{\mathrm{spec}},
\mathcal E_{\mathrm{ledger}},
\mathcal C_{\mathrm{budget}}
\bigr)
\]
为本文的首选最小实例簇，其中：
\begin{enumerate}
  \item \(\mathcal P_{\mathrm{LOST}}\) 取为观察者时空逻辑系统
  \[
  \mathrm{LOST}
  :=
  \bigl(
  P,\le,\Vdash,\{\mathrm{Upd}_\alpha\},J,I,\iota,
  X,\preceq_X,\mathcal U_X,\mathrm{Cov}_X,\lambda,\varrho,\{\mathrm{Com}_\mu\},
  \{C_K,\rho_i^K,\mathrm{supp}_K\}_{\varnothing\neq K\subseteq_{\mathrm{fin}} I}
  \bigr),
  \]
  其中观察者只是状态纤维索引，时间投影由
  \[
  \tau=\chi\circ\lambda
  \]
  给出；
  \item \(\mathcal X_{\mathrm{Fold}}\) 取为 Fold 读出的 restriction 逆系统
  \[
  \cdots \xrightarrow{\pi_{m+1,m}} X_m \xrightarrow{\pi_{m,m-1}} X_{m-1}\xrightarrow{}\cdots\xrightarrow{}X_1,
  \]
  其中
  \[
  X_m=\{w\in\{0,1\}^m:\ \text{不含子词 }11\},
  \]
  并由 \(\Fold_m:\Omega_m\to X_m\) 与层间 restriction 组织成逆系统，其无穷对象满足
  \[
  X_\infty\cong \varprojlim X_m;
  \]
  \item \(\mathcal U_{\mathrm{loc}}\) 取为前缀站点、局部覆盖、局部截面与重叠转移组成的局部可读对象系统；
  \item \(\mathcal R_{\mathrm{spec}}\) 取为有限状态核可计算的谱层，至少包括
  \[
  S_q(m),\qquad
  P(q),\qquad
  \zeta\text{-接口},\qquad
  \text{primitive 轨道谱},\qquad
  \mathcal L_q;
  \]
  \item \(\mathcal E_{\mathrm{ledger}}\) 取为残差外置层，包含 Fold 纤维残差、\(\mathrm{NULL}\)/失败证书、动态 prime-register 以及相位—账本双轴中的全不连通账本核；
  \item \(\mathcal C_{\mathrm{budget}}\) 取为资源预序层：对局部比较证书 \(c\) 定义预算向量
  \[
  \mathbf B(p,q;c)=(S,T,\Pi,L),
  \]
  再经单调次可加泛函 \(\Lambda\) 压成资源准距离
  \[
  d_{\mathrm{res}}(p,q):=\inf_c \Lambda(\mathbf B(p,q;c)).
  \]
\end{enumerate}
\end{definition}

定义 \ref{def:typed-address-biaxial-completion-preferred-minimal-instantiation} 的关键并不在于“把很多对象堆在一起”，而在于这六个分量之间都已经有正文给出的硬接口。
观察者时空状态模型在定义 \ref{def:logic-expansion-observer-spacetime-state-model} 及其后的因果相容链中已经成形；
restriction 逆系统与逆极限对象由命题 \ref{prop:recursive-addressing-refinement} 与定理 \ref{thm:recursive-addressing-dyadic-inverse-limit-stable-object} 给出严格对象化；
\(\mathrm{NULL}\) 三分解与 Čech-\(H^2\) 障碍由定理 \ref{thm:recursive-addressing-readout-separatedness-null}、定义 \ref{def:prefix-site-h2-gluing-obstruction} 与命题 \ref{prop:null-as-h2-obstruction} 已经对象化；
prime-register 的有限秩障碍与相位—账本短正合列分别由推论 \ref{cor:killo-no-finite-additive-register-linearization}、\S\ref{subsubsec:emergent-arithmetic-dynamic-prime-register-godel-dynamics} 与定理 \ref{thm:killo-localization-circle-dimension-prime-ledger-splitting} 给出；
而谱层则已由 \S\ref{subsubsec:frt-apparent-randomness-freezing-rough-statistics} 的 moment-kernel 与 projective transfer 接口焊接完毕。

\subsubsection{六层的不可删性}\label{subsubsec:typed-address-biaxial-completion-preferred-minimal-instantiation-minimality}
所谓“最小”，并不是说 \(\mathcal I_{\min}\) 的成分最少，而是说：若删去其中任一层，上一节的粗糙时空四元组就至少有一项失去对象层宿主。

\begin{proposition}[六层的不可删性]\label{prop:typed-address-biaxial-completion-preferred-minimal-instantiation-irreducible}
对首选最小实例簇 \(\mathcal I_{\min}\)，六个分量分别承担如下不可替代角色：
\begin{enumerate}
  \item 删去 \(\mathcal P_{\mathrm{LOST}}\)，则失去观察者纤维、更新、通信与事件定位，因而 \(\preceq\) 与 \(\tau\) 无法定义；
  \item 删去 \(\mathcal X_{\mathrm{Fold}}\)，则失去跨分辨率相容性与逆极限统一对象，因而“不同 \(m\) 只是同一本体的不同精细度”失去严格语义；
  \item 删去 \(\mathcal U_{\mathrm{loc}}\)，则失去局部可读—全局不可拼的语义宿主，因而 \(\Omega\) 只能退化为非对象化修辞；
  \item 删去 \(\mathcal R_{\mathrm{spec}}\)，则失去高阶碰撞、pressure、\(\zeta\) 与 continuous pressure 接口，因而 \(d_{\mathrm{res}}\) 中的谱预算 \(\Pi\) 无从量化；
  \item 删去 \(\mathcal E_{\mathrm{ledger}}\)，则失去纤维残差的外置承载，也失去乘法层与历史层的忠实保存；
  \item 删去 \(\mathcal C_{\mathrm{budget}}\)，则 \(\preceq,\tau,\Omega\) 虽可局部存在，但它们不再能被压成几何代价，因而“地形图”只剩定性图景。
\end{enumerate}
\end{proposition}

\begin{proof}
逐项检查上一节粗糙时空四元组的定义即可。
\(\preceq\) 与 \(\tau\) 由状态—事件—通信数据支撑；
\(d_{\mathrm{res}}\) 需要比较证书、资源向量与其标量化；
\(\Omega\) 则需要局部截面、重叠转移与跨尺度 residual 的对象层宿主。
若删去任一分量，则至少有一项从可审计对象退化为无宿主的非正式描述。
因此这六层在 \(\mathcal I_{\min}\) 中是不可删的。证毕。
\end{proof}

\subsubsection{从首选实例簇抽取粗糙时空骨架}\label{subsubsec:typed-address-biaxial-completion-preferred-minimal-instantiation-extraction}
现在可以把上一节的抽取命题在 \(\mathcal I_{\min}\) 上写成具体构造。

\begin{theorem}[首选最小实例簇的粗糙时空抽取]\label{thm:typed-address-biaxial-completion-preferred-minimal-instantiation-extraction}
由首选最小实例簇 \(\mathcal I_{\min}\) 可自然诱导一个粗糙时空四元组
\[
\mathfrak G_{\min}
:=
(\preceq_{\min},\tau_{\min},d_{\mathrm{res},\min},\Omega_{\min}),
\]
其构造如下：
\begin{enumerate}
  \item \textbf{因果前序：}
  令
  \[
  p\preceq_{\min} q
  \]
  当且仅当存在有限条由内部更新 \(\mathrm{Upd}_\alpha\) 与通信 \(\mathrm{Com}_\mu\) 组成的链，把 \(p\) 带到 \(q\)；
  \item \textbf{时间投影：}
  定义
  \[
  \tau_{\min}(p):=\chi(\lambda(p));
  \]
  \item \textbf{资源准距离：}
  对任意可比较状态对 \(p,q\)，令比较证书 \(c\) 同时记录共同支撑、同步深度、补账需求与谱分辨率，并形成预算向量
  \[
  \mathbf B(p,q;c)=\bigl(S,T,\Pi,L\bigr),
  \]
  然后定义
  \[
  d_{\mathrm{res},\min}(p,q):=\inf_c\Lambda(\mathbf B(p,q;c));
  \]
  \item \textbf{残差/曲率候选：}
  定义
  \[
  \Omega_{\min}=\Omega_{\mathrm{cov}}\boxplus \Omega_{\mathrm{scale}},
  \]
  其中 \(\Omega_{\mathrm{cov}}\) 来自覆盖重叠上的 Čech/holonomy residual，\(\Omega_{\mathrm{scale}}\) 来自跨分辨率图表上的不交换 residual。
\end{enumerate}
\end{theorem}

\begin{proof}
\(\preceq_{\min}\) 来自 \(\mathcal P_{\mathrm{LOST}}\) 中的更新与通信闭包；
\(\tau_{\min}\) 来自事件定位投影 \(\lambda\) 与时钟 \(\chi\) 的复合；
\(d_{\mathrm{res},\min}\) 来自 \(\mathcal C_{\mathrm{budget}}\) 对 \(\mathcal X_{\mathrm{Fold}}\)、\(\mathcal E_{\mathrm{ledger}}\) 与 \(\mathcal R_{\mathrm{spec}}\) 的联合标量化；
\(\Omega_{\min}\) 则由 \(\mathcal U_{\mathrm{loc}}\) 与 \(\mathcal X_{\mathrm{Fold}}\) 的覆盖 residual 与尺度 residual 共同给出。
各分量都在 \(\mathcal I_{\min}\) 的六层数据中具有明确对象层宿主，因此 \(\mathfrak G_{\min}\) 良定义。证毕。
\end{proof}

\subsubsection{这不是流形替代品，而是流形之前的一层}\label{subsubsec:typed-address-biaxial-completion-preferred-minimal-instantiation-before-manifold}
需要再次强调：\(\mathfrak G_{\min}\) 仍然不是现成的流形几何。
它更像一个由资源、粘接与尺度三者耦合形成的粗糙前几何接口。

\begin{remark}[为何它先于流形]\label{rem:typed-address-biaxial-completion-preferred-minimal-instantiation-before-manifold}
在 \(\mathfrak G_{\min}\) 中：
\begin{enumerate}
  \item \(\preceq_{\min}\) 是前序，不一定能由平滑光锥张量刻画；
  \item \(\tau_{\min}\) 是事件投影，不一定与同步深度同一化；
  \item \(d_{\mathrm{res},\min}\) 是扩展准伪度规，不一定对称，也可能取 \(+\infty\)；
  \item \(\Omega_{\min}\) 同时含覆盖方向 residual 与尺度方向 residual，不一定能被单一联络 \(1\)-形式吸收。
\end{enumerate}
因此，它不是“比广义相对论更一般的张量理论”，而是更靠前的一层：局部性、因果性、预算学与粘接失败先出现；流形与度规只是在这些结构进一步对称化、连续化与可积化之后才可能出现。
\end{remark}

\subsubsection{相位—账本双轴在首选实例簇中的位置}\label{subsubsec:typed-address-biaxial-completion-preferred-minimal-instantiation-phase-ledger}
上一节已经把“连通相位因子 + 全不连通账本因子”的双轴模板抽了出来。
在 \(\mathcal I_{\min}\) 里，这个双轴不是装饰，而正好对应两类完全不同的残差。

\begin{proposition}[\texorpdfstring{\(\mathcal I_{\min}\)}{I_min} 中的相位—账本双轴]\label{prop:typed-address-biaxial-completion-preferred-minimal-instantiation-phase-ledger}
在首选最小实例簇中，连通相位轴与全不连通账本轴分别承担：
\begin{enumerate}
  \item \textbf{相位轴：}承载局部连续读出、局部相容、闭路 holonomy 与相位型 residual；
  \item \textbf{账本轴：}承载折叠残差、dynamic prime-register 历史、可逆补账与前缀恢复代价。
\end{enumerate}
特别地，局部化对偶短正合列
\[
0\to \prod_{p\in S}\ZZ_p\to \widehat G_S\to \TT\to 0
\]
提供了这两条轴的最小模型：\(\TT\) 对应连通相位，\(\prod_{p\in S}\ZZ_p\) 对应全不连通素数账本核。
\end{proposition}

\begin{proof}
定理 \ref{thm:killo-localization-circle-dimension-prime-ledger-splitting} 已经证明：在该短正合列中，连通圆因子与全不连通 \(p\)-adic 核在结构上不可混同。
前者负责连续局部 glue 与 holonomy 语义，后者负责乘法层、前史层与残差层的外置承载。
它们共存于同一扩张对象中，却又不能互相替代；这正与 \(\mathcal I_{\min}\) 中 \(\mathcal U_{\mathrm{loc}}\) 与 \(\mathcal E_{\mathrm{ledger}}\) 的分工完全一致。证毕。
\end{proof}

命题 \ref{prop:typed-address-biaxial-completion-preferred-minimal-instantiation-phase-ledger} 在叙事上非常关键，因为它把两条此前容易混淆的主线彻底区分清楚：
相位不是账本的另一种写法，账本也不是相位的离散化近似；它们是同一扩张对象中的两条不同方向，一条负责连续 local glue，一条负责离散 external bookkeeping。

\subsubsection{为何它天然适合粗糙连续与表观随机}\label{subsubsec:typed-address-biaxial-completion-preferred-minimal-instantiation-roughness-randomness}
首选最小实例簇的价值，并不只是把已有章节拼起来，而是第一次把粗糙几何、表观随机与时空骨架收进同一个对象层容器。

\begin{remark}[粗糙可见量、表观随机与时空骨架的统一宿主]\label{rem:typed-address-biaxial-completion-preferred-minimal-instantiation-unified-host}
在 \(\mathcal I_{\min}\) 中：
\begin{enumerate}
  \item 粗糙可见量来自 \(\mathcal X_{\mathrm{Fold}}+\mathcal E_{\mathrm{ledger}}\) 的组合：Fold 负责把微态压到稳定层，ledger 则把未外置的残差再以逐尺度读出的方式注回可见层，于是得到 \S\ref{subsubsec:groupoid-zeckendorf-multiscale-rough-readout} 中那类“确定性但处处不可导”的可见轨道；
  \item 表观随机来自 \(\mathcal X_{\mathrm{Fold}}+\mathcal R_{\mathrm{spec}}\) 的组合：纤维多重度谱、高阶碰撞矩、压力与 mixed collision 会在有限分辨率下制造出“低阶像扩散，高阶快速冻结”的统计外观，但这些量仍然可由有限状态核与射影压力算子完全审计；
  \item 粗糙时空来自 \(\mathcal P_{\mathrm{LOST}}+\mathcal U_{\mathrm{loc}}+\mathcal C_{\mathrm{budget}}\) 的组合：状态、事件、区域、覆盖与预算学共同决定了“何时可比”“何处可比”“比较要花多少成本”“绕一圈会剩下什么”。
\end{enumerate}
因此，粗糙几何、表观随机与时空骨架并不是三套平行叙事，而是 \(\mathcal I_{\min}\) 的三种不同投影。
\end{remark}

\subsubsection{最小可审计实例的研究顺序}\label{subsubsec:typed-address-biaxial-completion-preferred-minimal-instantiation-order}
到这里，首选实例簇已经写出来了。
下一步不应立刻追求完整协变物理，而应先构造最小可审计版本。

\begin{remark}[最小研究顺序]\label{rem:typed-address-biaxial-completion-preferred-minimal-instantiation-order}
最自然的起点是固定一个最小窗口族
\[
m\in\{6,8,10,12\},
\]
并以 window-\(6\) 的精确纤维直方图
\[
|X_6|=21,
\qquad
2:8,\ 3:4,\ 4:9
\]
作为第一批锚点。
在这一锚点上，可逐项填入 \(\mathcal I_{\min}\) 的六层数据：先做 Fold 与可见壳，再做局部截面与 \(\mathrm{NULL}\)/失败证书，再做 mixed collision 与射影压力，最后把 dynamic prime-register 与预算向量挂进去。
因为 window-\(6\) 已经给出精确纤维直方图、规范体积密度、折叠缺陷与交换化秩等闭式对象，它正好适合充当第一批可复核样例。
\end{remark}

\subsubsection{小结：首选最小实例簇作为统一宿主}\label{subsubsec:typed-address-biaxial-completion-preferred-minimal-instantiation-summary}
本节完成的事可以压成一句话：本文最自然的物理型实例，并不是单独的符号动力系统、单独的 prime-register、单独的 Čech 覆盖或单独的资源函子，而是它们组成的复合实例簇 \(\mathcal I_{\min}\)；只有在这六层同时存在时，粗糙时空四元组
\[
(\preceq,\tau,d_{\mathrm{res}},\Omega)
\]
才真正拥有对象层宿主。

于是，前面新增的几条主线也就都可被读成同一实例簇的不同投影：
\begin{enumerate}
  \item \S\ref{subsubsec:groupoid-zeckendorf-multiscale-rough-readout} 是 \(\mathcal X_{\mathrm{Fold}}+\mathcal E_{\mathrm{ledger}}\) 的可见几何投影；
  \item \S\ref{subsubsec:frt-apparent-randomness-freezing-rough-statistics} 是 \(\mathcal X_{\mathrm{Fold}}+\mathcal R_{\mathrm{spec}}\) 的统计投影；
  \item \S\ref{subsubsec:emergent-arithmetic-dynamic-prime-register-godel-dynamics} 是 \(\mathcal E_{\mathrm{ledger}}\) 的动态乘法/历史投影；
  \item 上两节则是 \(\mathcal P_{\mathrm{LOST}}+\mathcal U_{\mathrm{loc}}+\mathcal C_{\mathrm{budget}}\) 的前几何投影。
\end{enumerate}
顺着这一点，后续各节最自然的任务不再是继续扩展抽象模板，而是把 \(\mathcal I_{\min}\) 真正落到一个最小可审计样例上。
第一步正是下一节的 window-\(6\) 样例：在那里，粗糙可见量、表观随机、boundary holonomy 与动态账本将第一次被压到同一个有限对象上。
